% This file is included by uet_thesis.tex.
% Replace blue placeholders and remove example text before submission.

\chapter{Literature Review and Related Work}
\label{chap:literature}
% =============================================================================

\section{Review Method}
\begin{guidance}
Briefly state how related work was identified and selected. Report the sources
searched, key search terms, publication period, and principal inclusion and
exclusion criteria. Follow departmental requirements for the minimum number and
type of sources.
\end{guidance}

\begin{table}[H]
\centering
\caption{Review method summary}
\label{tab:review-method}
\begin{tabularx}{\textwidth}{@{}p{0.27\textwidth}X@{}}
\toprule
\textbf{Item} & \textbf{Description} \\
\midrule
Sources searched & \placeholder{e.g., IEEE Xplore, ACM Digital Library, Scopus, official product sources} \\
Search terms & \placeholder{Main keywords and combinations} \\
Coverage period & \placeholder{Years covered and search date} \\
Selection criteria & \placeholder{Principal inclusion and exclusion criteria} \\
\bottomrule
\end{tabularx}
\end{table}

\section{Related Systems and Solutions \applicability{Development, Hybrid, Game}}
\begin{guidance}
Review relevant applications, systems, projects, or games. For each, explain the
problem addressed, important features, strengths, and limitations relevant to the
proposed project. Avoid promotional description.
\end{guidance}

\begin{table}[H]
\centering
\caption{Related systems and solutions}
\label{tab:related-systems}
\begin{tabularx}{\textwidth}{@{}p{0.20\textwidth}X X X@{}}
\toprule
\textbf{System and citation} & \textbf{Relevant features} & \textbf{Strengths} & \textbf{Limitations} \\
\midrule
\placeholder{System [ref]} & \placeholder{Features} & \placeholder{Strengths} & \placeholder{Limitations} \\
\placeholder{System [ref]} & \placeholder{Features} & \placeholder{Strengths} & \placeholder{Limitations} \\
\bottomrule
\end{tabularx}
\end{table}

\section{Comparative Analysis}
\begin{guidance}
Compare the most relevant systems or studies using criteria derived from the
project problem and objectives. Explain the important patterns shown by the table.
\end{guidance}

\begin{table}[H]
\centering
\caption{Comparative analysis}
\label{tab:comparative-analysis}
\begin{tabularx}{\textwidth}{@{}p{0.26\textwidth}X X X@{}}
\toprule
\textbf{System or study} & \textbf{Criterion 1} & \textbf{Criterion 2} & \textbf{Criterion 3} \\
\midrule
\placeholder{Name [ref]} & \placeholder{Finding} & \placeholder{Finding} & \placeholder{Finding} \\
\textbf{Proposed project} & \placeholder{Planned position} & \placeholder{Planned position} & \placeholder{Planned position} \\
\bottomrule
\end{tabularx}
\end{table}

\section{Gap Analysis and Proposed Contribution}
\begin{guidance}
State the specific unresolved gap established by the review and explain how the
project will address it. Keep the proposed contribution realistic and aligned with
the problem statement and objectives; do not claim novelty without evidence.
\end{guidance}

\placeholder{State the identified gap and the project's intended contribution.}

% =============================================================================
